\section{M6.4--M7 Parallel Blueprint Fragment}
\label{sec:m6_4_m7_parallel_blueprint_fragment}
\begingroup
\sloppy
\emergencystretch=3em

This fragment is the next handoff layer after the M6 constructor split.  Its
purpose is not to introduce another final theorem package.  The constructor
layer is already in place.  The remaining work is mathematical field
production: compact-support geometry, localized smoothness, artificial-face
geometry, boundary chart-change identification, and global integral
reconstruction.

Each entry below is organized by Lean module, record field, and theorem anchor.
The ``constructor status'' paragraph says what can already be used without
additional design.  The ``explicit field mathematics'' paragraph names the
field-level facts that are still mathematical work.  The ``parallel theorem
anchors'' paragraph is the intended worker queue.

\begin{definition}[Completed constructor layer]
  \label{anc:m6_4_completed_constructor_layer}
  \lean{Stokes.GlobalStokesData, Stokes.ProjectLocalGlobalStokesData,
    Stokes.PartitionReconstructionData, Stokes.BulkIntegralReconstructionData,
    Stokes.ExtDerivOnSupportData, Stokes.MixedGlobalStokesData,
    Stokes.SelectedMixedGlobalInput, Stokes.BoundaryGlobalConstructorData,
    Stokes.ArtificialFaceGeometryData, Stokes.BoundaryChartChangeFamilyData,
    Stokes.LocalizedInteriorPieces}
  \leanok
  \uses{anc:m6_constructor_dependency_order, anc:global_stokes_data}
  \emph{Lean modules:}
  the completed constructor modules are
  \texttt{Theorem}, \texttt{Reconstruction},
  \texttt{ReconstructionWrappers}, \texttt{IntegralReconstruction},
  \texttt{ExtDerivReconstruction}, \texttt{ExtDerivOnSupport},
  \texttt{MixedGlobalConstructor}, \texttt{MixedSelectedConstructor},
  \texttt{BoundaryGlobalConstructor}, \texttt{BoundaryChartChangePieces},
  \texttt{ArtificialFacePairing}, \texttt{ArtificialFaceGeometry}, and
  \texttt{LocalizedInteriorPieces}, all under \texttt{Stokes.Global}.

  \emph{Constructor status:}
  these modules already turn named local packages, cancellation packages,
  chart-change packages, and reconstruction packages into
  \texttt{GlobalStokesData}, \texttt{ProjectLocalGlobalStokesData}, or
  \texttt{MixedGlobalStokesData}.  The final theorems
  \texttt{globalStokes}, \texttt{projectLocalGlobalStokes},
  \texttt{mixedGlobalStokes}, and \texttt{selectedMixedGlobalStokes} are now
  bookkeeping theorem calls once their records are filled.

  \emph{Fields still mathematical:}
  do not add another global constructor to hide these fields.  Workers should
  instead prove the data needed for
  \texttt{globalBulkIntegral\_eq\_localBulkSum},
  \texttt{globalBoundaryIntegral\_eq\_boundaryPartitionSum},
  \texttt{interiorLocalStokes}, \texttt{boundaryLocalStokes},
  \texttt{interiorBoundaryCancellation},
  \texttt{chartChangeCancellation}, and the support-local
  \texttt{chartwiseExtDeriv\_eq\_global\_on} field.

  \emph{Handoff rule:}
  M6.4 workers should output one of the existing data packages, not a new
  final wrapper.  M7 workers should assemble those outputs through
  \texttt{SelectedMixedGlobalInput} or through
  \texttt{BulkIntegralReconstructionData.toMixedGlobalStokesData}.
\end{definition}

\begin{definition}[M6.4 active compact support to selected boxes]
  \label{anc:m6_4_active_compact_support_to_boxes}
  \lean{Stokes.ChartCompactImage, Stokes.ActiveChartCompactImages,
    Stokes.ChartCompactSupportData, Stokes.ActiveChartCompactSupportData,
    Stokes.CompactActiveBoxData, Stokes.CompactActiveExtendedBoxData}
  \leanok
  \uses{anc:m6_compact_active_boxes, def:selected_chart_box_cover,
    thm:finite_selected_box_cover_exists}
  \emph{Lean modules:}
  \texttt{ChartCompactImage}, \texttt{CompactSupportChartBox}, and
  \texttt{CompactActiveBoxes} under \texttt{Stokes.Global}.

  \emph{Constructor status:}
  \texttt{ChartCompactImage} and \texttt{ActiveChartCompactImages} already
  convert compact manifold control sets into compact coordinate images and
  selected coordinate boxes.  \texttt{ChartCompactSupportData} and
  \texttt{ActiveChartCompactSupportData} already carry the
  \texttt{tsupport} containment that converts those boxes into
  \texttt{interiorChartSelectedBox}.  \texttt{CompactActiveBoxData} and
  \texttt{CompactActiveExtendedBoxData} already expose lower corners, upper
  corners, support containment, target containment, overlap containment, and
  selected or extended box accessors.

  \emph{Explicit field mathematics:}
  the hard fields are still the compact coordinate support for each active
  chart, compactness of that support, the
  \texttt{tsupport\_subset\_coordSupport} proof, and the geometric shrinkage
  facts
  \[
    \operatorname{Icc}(a_i,b_i) \subseteq
      (\operatorname{extChartAt}\ I\ i).\operatorname{target},
    \qquad
    \operatorname{Icc}(a_i,b_i) \subseteq
      \operatorname{chartOverlap}(I,i,i).
  \]
  For extended boxes, workers must also supply an open neighborhood containing
  each box and the corresponding \texttt{ContDiffOn} proof.

  \emph{Parallel theorem anchors:}
  one worker can prove compact coordinate-image selection from a compact
  support set; one can prove target and overlap shrinkage for boxes chosen
  inside chart domains; one can produce the support-containment theorem for
  chart representatives; and one can lift selected boxes to extended boxes via
  localized smoothness.  Proposed Lean names:
  \plannedlean{Stokes.activeChartCompactSupportData_of_compactSupport}
  \plannedlean{Stokes.compactActiveBoxData_of_activeChartCompactSupport}
  \plannedlean{Stokes.compactActiveExtendedBoxes_of_smoothness}
\end{definition}

\begin{definition}[M6.4 partition localization and support control]
  \label{anc:m6_4_partition_localization_support_control}
  \lean{Stokes.FiniteActiveOnCompact, Stokes.SelectedBoxPartitionOfUnity,
    Stokes.LocalizedFormReconstructionFields, Stokes.LocalizedSupportControl,
    Stokes.LocalizedSmoothnessData, Stokes.ExtDerivOnSupportData}
  \leanok
  \uses{anc:m6_partition_sum_one, anc:m6_localized_support_smoothness,
    lem:partition_piece_support, lem:partition_sum_form}
  \emph{Lean modules:}
  \texttt{Stokes.Global.FiniteActive},
  \texttt{Stokes.Global.PartitionSumOne},
  \texttt{Stokes.Global.LocalizedSupport},
  \texttt{Stokes.Global.CoefficientBoxSupport},
  \texttt{Stokes.Global.LocalizedSmoothness},
  \texttt{Stokes.Global.ReconstructionWrappers}, and
  \texttt{Stokes.Global.ExtDerivOnSupport}.

  \emph{Constructor status:}
  \texttt{FiniteActiveOnCompact} already supplies the finite active set and
  the active-support membership lemmas.  \texttt{PartitionSumOne} already turns
  finite activity into the coefficient-sum-one facts used by
  \texttt{LocalizedFormReconstructionFields.ofCoeffSumEqOneOn}.
  \texttt{LocalizedSupportControl} already converts coefficient support in a
  selected box into selected-box support for the localized form
  \(\rho_i\omega\).  \texttt{LocalizedSmoothnessData} already packages the
  open-neighborhood smoothness field needed by extended boxes.

  \emph{Explicit field mathematics:}
  workers still have to prove that the coefficient family used in the
  reconstruction is exactly the selected partition coefficient, that each
  coefficient has the stated compact or boxed support, that
  \(\rho_i\omega\) is smooth in the selected chart neighborhood, and that the
  localized finite sum reconstructs \(\omega\) on the support set used by the
  global integrals.  The support-local exterior-derivative field
  \texttt{ExtDerivOnSupportData.chartwiseExtDeriv\_eq\_global\_on} is still a
  theorem, not a constructor obligation that should be guessed.

  \emph{Parallel theorem anchors:}
  a partition worker can prove the coefficient-sum-one and localized-form
  reconstruction on the compact control set; a support worker can build
  \texttt{LocalizedSupportControl} from coefficient support; a smoothness
  worker can prove \texttt{LocalizedSmoothnessData} from chartwise smoothness
  plus scalar smoothness; and an exterior-derivative worker can prove the
  support-local equality using
  \texttt{extDeriv\_transitionPullbackInChart\_localizedFormSum}.
  Proposed Lean names:
  \plannedlean{Stokes.selectedPartition_localizedFormReconstruction}
  \plannedlean{Stokes.localizedSupportControl_of_partitionCoeff}
  \plannedlean{Stokes.extDerivOnSupport_of_partitionSumOne}
\end{definition}

\begin{definition}[M6.4 interior local pieces and artificial faces]
  \label{anc:m6_4_interior_pieces_artificial_faces}
  \lean{Stokes.LocalizedInteriorPiece, Stokes.LocalizedInteriorPieces,
    Stokes.InteriorLocalStokesData, Stokes.ArtificialFaceGeometryData,
    Stokes.ArtificialFacePairingData, Stokes.ArtificialBoundaryCancellationData}
  \leanok
  \uses{anc:m6_localized_interior_pieces,
    anc:m6_artificial_face_pairing, thm:interior_chart_local_stokes}
  \emph{Lean modules:}
  \texttt{Stokes.Global.LocalizedInteriorPieces},
  \texttt{Stokes.Global.InteriorLocalStokes},
  \texttt{Stokes.Global.ArtificialFacePairing},
  \texttt{Stokes.Global.ArtificialFaceGeometry}, and
  \texttt{Stokes.Global.Cancellation}.

  \emph{Constructor status:}
  one \texttt{LocalizedInteriorPiece} already gives a localized form, selected
  box, extended box, \texttt{InteriorLocalStokesData}, and project-local
  equality.  A finite \texttt{LocalizedInteriorPieces} family already sums
  those local equalities.  \texttt{ArtificialFaceGeometryData} already forgets
  to \texttt{ArtificialFacePairingData} and then to
  \texttt{ArtificialBoundaryCancellationData}.

  \emph{Explicit field mathematics:}
  the remaining mathematical fields are the actual face labels for every
  interior piece, the map from each local artificial-boundary term to the
  recorded face sum, the pairing on flattened face indices, proof that the
  pairing preserves the active finite set, involutivity, equality of paired
  geometric faces, equality of paired unsigned face integrals, and opposite
  cubical signs.  If the chosen pairing has no fixed points, use
  \texttt{ArtificialFaceGeometryData.ofCoordinateFixedPointFreePairing};
  otherwise fixed-point zero terms must be proved explicitly.

  \emph{Parallel theorem anchors:}
  one worker can construct the \texttt{LocalizedInteriorPieces} family from
  support and smoothness packages; one can define the geometric face labels and
  unsigned terms; one can prove the opposite-face pairing; and one can prove
  that the artificial boundary term used by the local Stokes package equals
  the face sum consumed by the cancellation package.  Proposed Lean names:
  \plannedlean{Stokes.localizedInteriorPieces_of_activeBoxes}
  \plannedlean{Stokes.artificialFaceGeometryData_of_boxes}
  \plannedlean{Stokes.interiorBoundaryTerm_eq_artificialFaceSum}
\end{definition}

\begin{definition}[M6.4 boundary local pieces and chart-change fields]
  \label{anc:m6_4_boundary_pieces_chart_change_fields}
  \lean{Stokes.BoundaryProjectLocalPieces,
    Stokes.OrientedBoundaryProjectLocalPieces,
    Stokes.BoundaryChartChangePieceData,
    Stokes.BoundaryChartChangeFamilyData,
    Stokes.OrientedBoundaryChartChangeFamilyData,
    Stokes.BoundaryGlobalConstructorData}
  \leanok
  \uses{anc:m6_boundary_chart_change_pieces,
    anc:m6_boundary_piece_constructors, thm:boundary_chart_local_stokes}
  \emph{Lean modules:}
  \texttt{BoundaryPieces}, \texttt{BoundaryChartChangePieces}, and
  \texttt{BoundaryGlobalConstructor} under \texttt{Stokes.Global}, together
  with the boundary chart modules \texttt{LocalStokes},
  \texttt{SelectedBoxImageConstructor}, and \texttt{BoundaryPieceConvenience}.

  \emph{Constructor status:}
  boundary local Stokes can already be packaged either as
  \texttt{BoundaryProjectLocalPieces} or as
  \texttt{OrientedBoundaryProjectLocalPieces}.  Selected-target and
  extended-target \texttt{BoundaryChartChangeFamilyData} already produce the
  finite \texttt{chartChangeCancellation} field.  Boundary-only global
  constructors are already available for smoke tests, but the mixed theorem
  should use the same chart-change data through \texttt{SelectedMixedGlobalInput}.

  \emph{Explicit field mathematics:}
  boundary workers still need target chart choices, target lower and upper
  corners, target selected or extended boundary boxes, image data identifying
  source lower faces with target boundary boxes, the oriented change-of-
  variables package for every active boundary piece, and the pointwise field
  \texttt{boundaryPartitionTerm\_eq}.  Separately, the boundary reconstruction
  worker must identify the summed \texttt{boundaryPartitionTerm} with the
  chosen global boundary integral.

  \emph{Parallel theorem anchors:}
  one worker can build source boundary local pieces from selected image data;
  one can construct target selected boxes and local inverse/image data; one can
  prove oriented change of variables piecewise; and one can prove the
  \texttt{boundaryPartitionTerm\_eq} identification used by the chart-change
  family.  Proposed Lean names:
  \plannedlean{Stokes.boundaryPieces_of_selectedBoxImageData_family}
  \plannedlean{Stokes.boundaryChartChangeFamilyData_of_imageData}
  \plannedlean{Stokes.boundaryPartitionTerm_eq_targetBoundary}
\end{definition}

\begin{definition}[M6.5 global integral and exterior-derivative reconstruction]
  \label{anc:m6_5_global_integral_extderiv_reconstruction}
  \lean{Stokes.BulkIntegralReconstructionData,
    Stokes.BulkIntegralReconstructionData.BoundaryPartitionFields,
    Stokes.PartitionReconstructionData,
    Stokes.GlobalStokesReconstructionFields,
    Stokes.ExtDerivPartitionReconstructionData,
    Stokes.ExtDerivOnSupportData}
  \leanok
  \uses{anc:m6_reconstruction_constructors,
    thm:global_bulk_eq_local_sum, thm:global_boundary_eq_local_sum,
    lem:extderiv_finite_sum}
  \emph{Lean modules:}
  \texttt{Stokes.Global.IntegralReconstruction},
  \texttt{Stokes.Global.Reconstruction},
  \texttt{Stokes.Global.ReconstructionWrappers},
  \texttt{Stokes.Global.ExtDerivReconstruction}, and
  \texttt{Stokes.Global.ExtDerivOnSupport}.

  \emph{Constructor status:}
  bulk reconstruction, boundary reconstruction, and exterior-derivative
  reconstruction have been split into reusable records.  A worker may first
  prove only \texttt{BulkIntegralReconstructionData}; another worker may add
  \texttt{BoundaryPartitionFields}; another may add
  \texttt{ExtDerivOnSupportData}.  These records already convert back to
  \texttt{PartitionReconstructionData} and to the final
  \texttt{GlobalStokesData} or \texttt{MixedGlobalStokesData} shapes.

  \emph{Explicit field mathematics:}
  the project still needs actual definitions or comparisons for the represented
  global bulk and boundary integrals.  The bulk field must prove finite
  additivity over localized pieces and the equality with the local
  project-integral terms.  The boundary field must prove finite additivity over
  the boundary partition and compatibility with the boundary chart-change
  target terms.  The exterior-derivative field must prove that the coordinate
  representative of \(\dd(\sum_i \rho_i\omega)\) equals the coordinate
  representative of \(\dd\omega\) on the support on which the bulk integral is
  reconstructed.

  \emph{Parallel theorem anchors:}
  one worker can implement the bulk integral definition and finite-additivity
  theorem; one can implement the boundary integral definition and
  finite-additivity theorem; one can prove the \texttt{extDeriv} support-local
  reconstruction; and one can write the comparison wrappers from these
  definitions to the project-local terms.  Proposed Lean names:
  \plannedlean{Stokes.globalBulkIntegral_eq_selectedLocalBulkSum}
  \plannedlean{Stokes.globalBoundaryIntegral_eq_boundaryPartitionSum}
  \plannedlean{Stokes.chartwiseExtDeriv_sum_eq_global_on}
\end{definition}

\begin{theorem}[M7 selected mixed global assembly]
  \label{anc:m7_selected_mixed_global_assembly}
  \lean{Stokes.SelectedMixedGlobalInput,
    Stokes.SelectedMixedGlobalInput.toMixedGlobalStokesData,
    Stokes.SelectedMixedGlobalInput.toGlobalStokesData,
    Stokes.SelectedMixedGlobalInput.stokes,
    Stokes.selectedMixedGlobalStokes, Stokes.globalStokes}
  \leanok
  \uses{anc:m6_4_active_compact_support_to_boxes,
    anc:m6_4_partition_localization_support_control,
    anc:m6_4_interior_pieces_artificial_faces,
    anc:m6_4_boundary_pieces_chart_change_fields,
    anc:m6_5_global_integral_extderiv_reconstruction}
  \emph{Lean module:} \texttt{Stokes.Global.MixedSelectedConstructor}.

  \emph{Constructor status:}
  \texttt{SelectedMixedGlobalInput} is the M7 parent-agent target.  Once its
  fields are supplied, \texttt{selectedMixedGlobalStokes} gives the final
  equality recorded by the reconstruction package.  The fields
  \texttt{selectedPartition\_active},
  \texttt{artificialCancellation\_active},
  \texttt{artificialCancellation\_pieces},
  \texttt{artificialCancellation\_term},
  \texttt{chartChange\_active}, \texttt{chartChange\_pieces},
  \texttt{chartChange\_oldTerm}, and \texttt{chartChange\_newTerm} are
  alignment fields; they should be solved by choosing common indices and common
  term functions, not by new analysis.

  \emph{Explicit field mathematics:}
  the only remaining non-bookkeeping inputs are the selected partition, the
  reconstruction package, the interior local package, the boundary local
  package, artificial-boundary cancellation, and chart-change cancellation.
  These are exactly the outputs of the M6.4 and M6.5 entries above.

  \emph{M7 theorem anchors:}
  the compact-support theorem should instantiate
  \texttt{SelectedMixedGlobalInput} from a compact support set, a selected
  partition of unity, local boxes, reconstruction data, and cancellation data.
  The compact-manifold theorem should reduce to compact support.  A later
  mathlib-facing theorem should compare the project-local integral API to any
  canonical manifold-form integral exposed upstream.  Proposed Lean names:
  \plannedlean{Stokes.globalStokes_compactSupport_from_input}
  \plannedlean{Stokes.globalStokes_compactSupport}
  \plannedlean{Stokes.globalStokes_compactManifold}
\end{theorem}

\begin{theorem}[Parallel theorem queue for M6.4--M7]
  \label{anc:m6_4_m7_parallel_theorem_queue}
  \uses{anc:m6_4_completed_constructor_layer,
    anc:m6_4_active_compact_support_to_boxes,
    anc:m6_4_partition_localization_support_control,
    anc:m6_4_interior_pieces_artificial_faces,
    anc:m6_4_boundary_pieces_chart_change_fields,
    anc:m6_5_global_integral_extderiv_reconstruction,
    anc:m7_selected_mixed_global_assembly}
  The next parallel pass should avoid record redesign.  The natural independent
  work units are:
  \begin{enumerate}
    \item compact support to active selected boxes;
    \item selected partition coefficient support and sum-one reconstruction;
    \item localized smoothness and support-local exterior derivative
      reconstruction;
    \item localized interior local Stokes packages;
    \item artificial-face geometry and cancellation;
    \item boundary selected-box image data and oriented chart change;
    \item global bulk reconstruction;
    \item global boundary reconstruction;
    \item final \texttt{SelectedMixedGlobalInput} assembly.
  \end{enumerate}
  Items 1--3 can run in parallel with the boundary chart-change work, provided
  all workers agree on the active chart labels and coefficient family.  Items
  4 and 5 can run as soon as support and smoothness packages exist.  Items 7
  and 8 should consume the same local term functions used by items 4--6, so the
  final M7 assembly is mostly field alignment.
\end{theorem}

\endgroup
